\scanpage{46}
catégoriques généraux --- ].

Donc en ce cas, l'invariant suffit à réconstituer \deleted{\unsure{toute|tout}} la \unsure{structure|situation} $(X,X',s)$.

Dans le cas où $X$ est normal, intègre, et où $i$ n'\uline{est pas identiquement égal} \uline{à $0$ ou $1$}, comme on \unsure{connaîtra|connaît} \unsure{le|la} \unsure{revê|couv.|revêtement}\deleted{\unsure{t|ts}} non ramifié de $X$ sur l'ouvert \inserted{$U$} où $i$ et $i-1$ inversibles $[\unsure{D_{i(i-1)}|D_{i,i-1}|S_{i(i-1)}|S_{j(j-1)}}]$, on \unsure{le|la} \unsure{connaîtra|connaît} partout, d'où reconstitution de $(X,X')$ ; d'autre part, on connaît $s$ sur $U$, donc ($X$ normal) partout.

Donc ici encore, \uline{$i$ permet de tt réconstruire}. Ici cependant, \unsure{se} pose la question : quels $i$ conviennent ??

Il faut exprimer (i) qu'une certaine extension de corps -- de degré $3$ -- du corps des fonctions de $X$ est non ramifié sur $X$ (ii) Qu'une certaine application rationnelle de $X$ qui s'en déduit est partout définie ----
% archive sheet number: 45
